% ============================================================
% SNIPPET PARA CAPÍTULO 5 — Sección de validación del gateway
% Insertar después de la descripción del ciclo watch / systemd
% ============================================================

\subsection{Validación de laboratorio y verificación en campo}
\label{subsec:validacion_gateway_campo}

La descripción de la arquitectura en las secciones anteriores cubre el diseño
declarado; esta sección documenta la verificación empírica de que la pasarela
se comporta conforme a ese diseño en condiciones reales de campo. La
verificación se articuló en dos planos complementarios: pruebas de dry-run
en Windows para confirmar la normalización y consistencia de la salida del
motor, y pruebas sobre el nodo Raspberry Pi para confirmar el comportamiento
del canal \ac{MQTT}/\ac{TLS} y la resiliencia del spool ante fallos del broker.

La suite cubrió las sesiones de ordeño completo del período posterior a la
intervención (11--28 de mayo de 2026, 33 sesiones con correlación efectiva en
17 visitas) y agrupó seis tipos de verificación. Las pruebas se dividieron en
dos grupos: las que dependen del contenido de cada sesión (esquema, idempotencia
y consistencia de $\eta$) se ejecutaron sobre las 33 sesiones mediante dry-run
en Windows; las que verifican propiedades del sistema (seguridad del canal,
resiliencia ante fallos y recepción confirmada) se ejecutaron sobre una muestra
representativa de 10 visitas del tramo 11--20 de mayo en la Raspberry Pi, más
una verificación spot sobre el nodo en operación (29/05/2026).

La Tabla~\ref{tab:gateway_validation_suite} resume los resultados.

\begin{table}[H]
\centering
\caption{Resultados de la suite de validación de la pasarela perimetral
         (período post-intervención, 17 visitas, 33 sesiones, 11--28 de mayo de 2026).}
\label{tab:gateway_validation_suite}
\footnotesize
\renewcommand{\arraystretch}{1.22}
\begin{tabularx}{\linewidth}{|>{\raggedright\arraybackslash}p{3.2cm}|>{\raggedright\arraybackslash}X|>{\centering\arraybackslash}p{2.0cm}|>{\raggedright\arraybackslash}p{3.8cm}|}
\hline
\rowcolor{gray!20}
\TableHeaderFirst{Prueba} & \TableHeader{Alcance} & \TableHeader{Resultado} & \TableHeader{Observación} \\ \hline

Esquema \acs{JSON} &
Estructura y tipos de campo de cada lote publicado &
\textbf{33/33} \acs{PASS}$^*$ &
Verificado por dry-run sobre cada sesión (Windows). \\ \hline

Idempotencia &
Reproducibilidad de la salida entre corridas aisladas &
\textbf{33/33} \acs{PASS} &
Verificado por dry-run sobre cada sesión; \acs{MD5} idéntico en directorios aislados. \\ \hline

Consistencia $\eta$ &
Exactitud de la métrica publicada vs.\ métrica del motor &
\textbf{33/33} \acs{PASS} &
$\eta_{\text{motor}} = \eta_{\text{gateway}}$ en las 33 sesiones
($\bar{\eta}_{\text{post}} = 26{,}25\%$; rango: $8{,}3\%$--$100{,}0\%$). \\ \hline

\acs{TLS} Handshake &
Cifrado y versión mínima del canal \acs{MQTT} &
\textbf{11/11} \acs{PASS} &
Verificado en muestra representativa de visitas en Pi; \acs{TLS} 1.2 rechazado,
1.3 aceptado. \\ \hline

Resiliencia del spool &
Continuidad ante caída del broker \acs{MQTT} &
\textbf{11/11} \acs{PASS} &
Lote completo drenado tras recuperación del broker; spool residual $= 0$. \\ \hline

Suscripción \acs{MQTT} &
Recepción confirmada por suscriptor externo &
\textbf{11/11} \acs{PASS} &
Mensajes verificados con \texttt{mosquitto\_sub} activo post-recuperación. \\ \hline

\end{tabularx}
\begin{flushleft}
\footnotesize $^*$El tipo \texttt{field\_validation\_summary} no se genera en sesiones
sin validación presencial simultánea, por lo que su ausencia no constituye un
error de normalización. Las pruebas de canal seguro y resiliencia son propiedades
del sistema; su verificación sobre muestra representativa de 10 visitas es
metodológicamente suficiente porque no dependen del contenido particular de cada
sesión.
\end{flushleft}
\end{table}

La prueba de resiliencia merece descripción adicional porque expone el
comportamiento de extremo a extremo del spool. La secuencia de verificación
fue: (i)~detener el broker \texttt{mosquitto}, (ii)~ejecutar el gateway para
que espole el lote en disco, (iii)~reiniciar el broker, (iv)~activar un
suscriptor externo e invocar el drenado, y (v)~confirmar que el suscriptor
recibe exactamente los mensajes esperados y que el spool queda vacío. El
Listado~\ref{lst:resiliencia_salida} muestra la salida representativa de esa
secuencia sobre una sesión de 13 mensajes.

\begin{lstlisting}[
  language=bash,
  caption={Salida del test de resiliencia en \texttt{gateway-esmeralda} (29/05/2026).},
  label=lst:resiliencia_salida,
  basicstyle=\ttfamily\footnotesize,
  breaklines=true,
  frame=single
]
--- Mosquitto detenido ---
13 /tmp/spool/TOMA_PM__1PM__Captura_20260518_130005.jsonl
--- Mosquitto reiniciado ---
drained=13  failed=0
--- Resultados ---
Mensajes recibidos : 13
Spool residual     : 0
[PASS] Resiliencia OK
\end{lstlisting}

Durante esta verificación se identificó que la entrega de mensajes con
\acs{QoS}$\geq 1$ requería que el bucle de red del cliente \texttt{paho-mqtt}
permaneciera activo hasta la confirmación del broker. Se corrigió
\texttt{publish\_batch} añadiendo \texttt{loop\_start()} antes de publicar y
\texttt{loop\_stop()} tras confirmar la entrega, garantizando que cada lote
se transmite completamente antes de cerrar la conexión sin alterar el
comportamiento del spool ni el protocolo de reintentos.

Con estas pruebas se cerró la verificación del componente de transporte. El
contraste estadístico de $\eta$ pre/post intervención se presenta en el
Capítulo~\ref{chapter:chapter06}.
